\documentclass[12pt]{article}
\usepackage[utf8]{inputenc}
\usepackage[T1]{fontenc}
\usepackage{hyperref}
\title{Yonaguni Language Resources for K-12 Curriculum Development}
\author{Compiled by Claude Sonnet 5.5 (claude-sonnet-5-5) for Andr\'as Kornai}
\date{}
\begin{document}
\maketitle

\section{Introduction}
Yonaguni (Dunan-munui) is a Japonic language of Yonaguni island, the westernmost island of Japan. UNESCO classes it as severely endangered. Speaker estimates conflict: about 400 on older census-based figures, and fewer than 100 on the island in a 2023 paper, which also says that natural transmission between generations stopped in the mid-1960s. The island has no senior high school, so children leave at 15. The search found a few school and community projects and one dictionary, but no full curriculum materials.

\section{Primary Educational Resources}
\begin{itemize}
\item A 2025 conference study by Sachiyo Fujita-Round describes a school and community programme on Yonaguni. Pupils made a bilingual Japanese/Dunan film and a Dunan word-chain game video: \url{https://nomadit.co.uk/conference/eajs2026/paper/96786}
\item A panel slide set from the JALT conference, Nov. 2025: \url{https://events.jalt.org/event/37/contributions/1273/attachments/221/490/2025Nov02_JALT_PANEL_Sachiyo.pdf}
\end{itemize}
I found no published primers or textbooks.

\section{Linguistic and Reference Materials}
\begin{itemize}
\item Nae Ikema's Yonaguni dictionary (2003), listed in the Endangered Languages Project: \url{https://www.endangeredlanguages.com/node/103868}
\item Endangered Languages Project page for Dunan: \url{https://endangeredlanguages.com/lang/3170}
\item Resources listing for Dunan: \url{https://www.endangeredlanguages.com/node/38730/resources}
\item Heinrich, paper on Ryukyuan languages (2023): \url{https://iris.unive.it/retrieve/a54d8141-80f9-4edb-8a4e-4a75025aeda6/Harararude.pdf}
\item NINJAL Ryukyuan revitalization summary: \url{https://www2.ninjal.ac.jp/ael/files/resume/2_1_YamadaMasahiro.pdf}
\item Writing and orthography overview: \url{https://www.omniglot.com/writing/yonaguni.htm}
\end{itemize}
A 2025 overview says no full-scale dictionary has been published. Only bilingual Dunan--Japanese word lists exist.

\section{Religious and Community Resources}
None found.

\section{Digital and Audio Resources}
\begin{itemize}
\item A video of Nae Ikema speaking Yonaguni about her dictionary appears in the Endangered Languages Project record above.
\item Overview article: \url{https://thelanguagecloset.com/2025/09/20/the-language-in-japans-westernmost-island/}
\end{itemize}

\section{Curriculum Development Context}
Transmission in families has stopped, so any curriculum would be a second-language or immersion curriculum led by the few remaining speakers. The school-based video workshops are the only school model found. A conference panel notes that Ryukyuan revitalization lacks long-term language planning: \url{https://nomadit.co.uk/conference/eajs2026/p/19095}. The Omniglot page mentions a language-nest immersion plan, which I could not verify further.
\end{document}